% physics_3d_mirror.tex
% 3d mirror symmetry, CY_2 quantum vertex chiral groups, and Coulomb branches

\documentclass[11pt]{article}
\usepackage[localtheorems]{raeez-math-template}


% Theorem environments
\newtheorem{theorem}{Theorem}[section]
\newtheorem{proposition}[theorem]{Proposition}
\newtheorem{lemma}[theorem]{Lemma}
\newtheorem{corollary}[theorem]{Corollary}
\newtheorem{conjecture}[theorem]{Conjecture}
\theoremstyle{definition}
\newtheorem{definition}[theorem]{Definition}
\newtheorem{example}[theorem]{Example}
\newtheorem{construction}[theorem]{Construction}
\theoremstyle{remark}
\newtheorem{remark}[theorem]{Remark}
\newtheorem{warning}[theorem]{Warning}

% Macros (matching main monograph)
\newcommand{\CY}{\mathrm{CY}}
\newcommand{\HH}{\mathrm{HH}}
\newcommand{\Rep}{\mathrm{Rep}}
\newcommand{\Hom}{\mathrm{Hom}}
\newcommand{\End}{\mathrm{End}}
\newcommand{\Sym}{\mathrm{Sym}}
\newcommand{\Spec}{\mathrm{Spec}}
\newcommand{\Mod}{\mathrm{Mod}}
\newcommand{\Coh}{\mathrm{Coh}}
\newcommand{\Perf}{\mathrm{Perf}}
\newcommand{\Fuk}{\mathrm{Fuk}}
\newcommand{\Bun}{\mathrm{Bun}}
\newcommand{\Higgs}{\mathrm{Higgs}}
\newcommand{\Tr}{\mathrm{Tr}}
\newcommand{\id}{\mathrm{id}}
\newcommand{\op}{\mathrm{op}}
\newcommand{\Eone}{E_1}
\newcommand{\Etwo}{E_2}
\newcommand{\cA}{\mathcal{A}}
\newcommand{\cC}{\mathcal{C}}
\newcommand{\cM}{\mathcal{M}}
\newcommand{\cN}{\mathcal{N}}
\newcommand{\cO}{\mathcal{O}}
\newcommand{\frakg}{\mathfrak{g}}
\newcommand{\frakn}{\mathfrak{n}}
\renewcommand{\hbar}{\hslash}
\newcommand{\CoHA}{\mathrm{CoHA}}

\title{3d Mirror Symmetry, $\CY_2$ Quantum Vertex Chiral Groups,\\
and Coulomb Branches}
\author{Raeez Lorgat}
\date{April 2026}

\begin{document}
\maketitle

\begin{abstract}
We formulate the $\CY_2$ root--BPS dictionary behind 3d $\cN=4$ mirror symmetry. A 3d $\cN=4$ gauge theory $T$ has two hyperk\"ahler branches of vacua, $M_H(T)$ and $M_C(T)$, exchanged by mirror symmetry:
\[
 M_H(T) \cong M_C(T^!), \qquad M_C(T) \cong M_H(T^!).
\]
On constructed branches the $\CY_2$ functor produces $\Etwo$-chiral algebras. The Higgs side is controlled by quiver-variety CoHA and Maulik--Okounkov stable envelopes; the Coulomb side is controlled by the Braverman--Finkelberg--Nakajima convolution algebra and its quantization. Mirror symmetry predicts equality of these $\Etwo$ data after the derived-category and $R$-matrix identifications have been supplied.
\end{abstract}

\tableofcontents

%% =====================================================================
\section{3d $\cN=4$ theories: Higgs and Coulomb branches}
\label{sec:3d-n4}
%% =====================================================================

\subsection{The two branches of vacua}

A 3d $\cN=4$ gauge theory $T$ is specified by a compact gauge group $G$ and a quaternionic representation $R$. The moduli space of supersymmetric vacua splits into two branches:

\begin{enumerate}[label=(\roman*)]
 \item The \emph{Higgs branch} $M_H(T)$, parametrized by the vacuum expectation values of hypermultiplet scalars. Classically,
 \[
 M_H(T) = T^*R \mathbin{/\!\!/\!\!/} G = \mu_{\mathbb{R}}^{-1}(0) \cap \mu_{\mathbb{C}}^{-1}(0) / G,
 \]
 the hyperk\"ahler quotient, where $\mu_{\mathbb{R}}, \mu_{\mathbb{C}}$ are the real and complex moment maps for $G$ acting on $T^*R$.

 \item The \emph{Coulomb branch} $M_C(T)$, parametrized by the vacuum expectation values of vector multiplet scalars and dual photons. Classically, for abelian $G = U(1)^r$, $M_C \cong (\mathbb{R}^3 \times S^1)^r$, but quantum corrections drastically modify the geometry for non-abelian $G$.
\end{enumerate}

\begin{theorem}[Hyperk\"ahler property]
\label{thm:hk-branches}
Both $M_H(T)$ and $M_C(T)$ carry hyperk\"ahler metrics. In particular:
\begin{enumerate}[label=(\alph*)]
 \item $M_H(T)$ is hyperk\"ahler by the hyperk\"ahler quotient construction (Hitchin--Karlhede--Lindstr\"om--Ro\v{c}ek).
 \item $M_C(T)$ is hyperk\"ahler by the analysis of the low-energy effective action on the Coulomb branch (Seiberg--Witten for 3d).
\end{enumerate}
A hyperk\"ahler manifold of complex dimension $2n$ (real dimension $4n$) is Ricci-flat K\"ahler with holonomy $\subseteq \mathrm{Sp}(n)$. In particular it is Calabi--Yau. When $\dim_{\mathbb{C}} = 2$ (i.e.\ $n = 1$, real dimension 4), the holonomy condition is $\mathrm{Sp}(1) \cong \mathrm{SU}(2)$, so the manifold is $\CY_2$ (a K3 surface or a deformation thereof, in the compact case; or a hyperk\"ahler resolution of a singularity, in the typical gauge-theory case).
\end{theorem}

\begin{remark}[The $\CY_2$ condition]
\label{rem:cy2}
In the conventions of this monograph, ``$\CY_d$'' means the category $\cC$ (e.g., $D^b(\Coh(X))$ or $\Fuk(X)$) has a Serre functor $S \cong [d]$, i.e., the non-degenerate trace $\Tr \colon \HH_d(\cC) \to k$ lives in degree $d$. For a hyperk\"ahler surface $X$ of complex dimension 2, both $D^b(\Coh(X))$ and $\Fuk(X)$ are $\CY_2$ categories. This is the dimension needed for the $\mathbb{S}^2$-framing of Hochschild homology to provide an $\Etwo$-algebra structure (Chapter~\ref*{ch:cy-to-chiral} of the main text).
\end{remark}

\subsection{3d mirror symmetry}

\begin{definition}[3d mirror symmetry]
\label{def:3d-mirror}
Two 3d $\cN=4$ theories $T$ and $T^!$ are \emph{3d mirrors} if there is an equivalence of their moduli of vacua that exchanges the two branches:
\[
 M_H(T) \cong M_C(T^!), \qquad M_C(T) \cong M_H(T^!).
\]
This equivalence is an isomorphism of hyperk\"ahler manifolds (up to a specific rotation of the hyperk\"ahler structure, swapping complex structures $I$ and $J$).
\end{definition}

This duality was discovered by Intriligator--Seiberg (1996) and developed by Hanany--Witten via brane constructions. The prototypical example:

\begin{example}[ADHM mirror]
\label{ex:adhm}
Let $T$ be 3d $\cN=4$ $U(k)$ gauge theory with $N$ fundamental hypermultiplets. Then:
\[
 M_H(T) = T^*\mathrm{Gr}(k, N)
\]
is the cotangent bundle of the Grassmannian (a Nakajima quiver variety of type $A_1$ with framing $N$), and the mirror $T^!$ has $M_C(T^!) = T^*\mathrm{Gr}(k, N)$ as its Coulomb branch. For $k=1$: $M_H = T^*\mathbb{CP}^{N-1}$ and $M_C = \mathbb{C}^2/\mathbb{Z}_N$ (the $A_{N-1}$ singularity), exchanged under mirror symmetry.
\end{example}

\begin{example}[Self-mirror: $T^*\mathbb{CP}^1$]
\label{ex:self-mirror}
The $A_1$ singularity $\mathbb{C}^2/\mathbb{Z}_2$ and its resolution $T^*\mathbb{CP}^1$ are exchanged under mirror symmetry with the FI/mass parameter exchange. In the language of Nakajima quiver varieties, this is the self-mirror property of the Jordan quiver with framing 2.
\end{example}

%% =====================================================================
\section{The Higgs branch: hyperk\"ahler quotients and quiver varieties}
\label{sec:higgs-branch}
%% =====================================================================

\subsection{Classical description}

For a 3d $\cN=4$ theory with gauge group $G$ and matter representation $R$ (viewed as a complex representation of $G_{\mathbb{C}}$), the Higgs branch is the hyperk\"ahler quotient:
\[
 M_H = T^*R \mathbin{/\!\!/\!\!/} G.
\]
In complex-algebraic terms (in complex structure $I$), this is the GIT quotient
\[
 M_H = \mu_{\mathbb{C}}^{-1}(0) /\!/ G_{\mathbb{C}} = \Spec\left(\mathbb{C}[\mu_{\mathbb{C}}^{-1}(0)]^{G_{\mathbb{C}}}\right),
\]
where $\mu_{\mathbb{C}} \colon T^*R \to \frakg^*$ is the complex moment map.

\begin{remark}[Non-renormalization]
\label{rem:non-renorm}
A fundamental result of 3d $\cN=4$ physics: the Higgs branch receives \emph{no quantum corrections}. The classical hyperk\"ahler quotient description is exact. This is in sharp contrast to the Coulomb branch, which is heavily quantum-corrected.
\end{remark}

\subsection{Quiver gauge theories and Nakajima quiver varieties}

The most important class of 3d $\cN=4$ theories are \emph{quiver gauge theories}, specified by a quiver $Q = (Q_0, Q_1)$ with:
\begin{itemize}
 \item A dimension vector $\mathbf{v} = (v_i)_{i \in Q_0}$: the gauge group is $G = \prod_{i \in Q_0} U(v_i)$.
 \item A framing vector $\mathbf{w} = (w_i)_{i \in Q_0}$: the matter content includes $w_i$ fundamental hypermultiplets at node $i$.
\end{itemize}

The Higgs branch is then the \emph{Nakajima quiver variety}:
\[
 M_H = \mathfrak{M}(\mathbf{v}, \mathbf{w}) = \left\{(B, i, j) \;\middle|\; \mu_{\mathbb{C}}(B,i,j) = 0, \text{ stability}\right\} / G_{\mathbb{C}},
\]
where $B = (B_a)_{a \in Q_1}$ are the maps along arrows, $i = (i_k)$ and $j = (j_k)$ are the framing maps, and $\mu_{\mathbb{C}} = \sum_{a \in Q_1} [B_a, B_{a^*}] + ij$ is the complex moment map.

\begin{theorem}[Nakajima]
\label{thm:nakajima}
The quiver variety $\mathfrak{M}(\mathbf{v}, \mathbf{w})$ is a smooth (for generic stability) hyperk\"ahler manifold of complex dimension $2\sum_i v_i w_i - 2\sum_i v_i^2 + 2\sum_{a \in Q_1} v_{s(a)} v_{t(a)}$. It carries a natural action of the ``flavor symmetry'' group $\prod_i GL(w_i)$.
\end{theorem}

%% =====================================================================
\section{The Coulomb branch: BFN construction}
\label{sec:coulomb-bfn}
%% =====================================================================

\subsection{The problem of defining Coulomb branches}

Unlike the Higgs branch, the Coulomb branch has no simple classical description for non-abelian $G$: quantum corrections (monopole operators, instanton effects) are essential. Before Braverman--Finkelberg--Nakajima, it was defined only indirectly (as the Higgs branch of the mirror).

\subsection{The BFN definition}

Braverman, Finkelberg, and Nakajima (2016--2019) gave a rigorous definition using a convolution algebra in the affine Grassmannian.

\begin{definition}[Coulomb branch algebra (BFN)]
\label{def:bfn}
Let $G$ be a complex reductive group and $R$ a representation of $G$. Define:
\begin{enumerate}[label=(\roman*)]
 \item The \emph{affine Grassmannian} $\mathrm{Gr}_G = G(\!(t)\!) / G[\![t]\!]$.
 \item The \emph{space of triples} $\mathcal{R} = \{(g, s) \mid g \in \mathrm{Gr}_G,\; s \in R[\![t]\!],\; g \cdot s \in R[\![t]\!]\}$ (the space of pairs of a lattice in $\mathrm{Gr}_G$ and a section of $R$ compatible with the lattice change).
 \item The \emph{BFN algebra}
 \[
 \mathcal{A}_{G,R} = H^{G[\![t]\!]}_\bullet(\mathcal{R})
 \]
 is the $G[\![t]\!]$-equivariant Borel--Moore homology of $\mathcal{R}$, equipped with a convolution product.
\end{enumerate}
The \emph{Coulomb branch} is
\[
 M_C(G, R) = \Spec(\mathcal{A}_{G,R}).
\]
\end{definition}

\begin{theorem}[BFN]
\label{thm:bfn}
The algebra $\mathcal{A}_{G,R}$ is a commutative, finitely generated $\mathbb{C}$-algebra. The variety $M_C(G, R)$ is a normal, affine, conical symplectic singularity (in the sense of Beauville). It admits a canonical Poisson structure and, upon quantization (replacing equivariant homology by equivariant $K$-theory), a canonical non-commutative deformation.
\end{theorem}

\begin{remark}[Quantized Coulomb branch]
\label{rem:quantized-coulomb}
The quantization of $\mathcal{A}_{G,R}$ is the \emph{quantized Coulomb branch algebra}
\[
 \mathcal{A}_{G,R}^\hbar =
 H^{G[\![t]\!] \rtimes \mathbb{C}^\times}_\bullet(\mathcal{R}),
\]
where the extra $\mathbb{C}^\times$ rotates the loop parameter and supplies $\hbar$. This is a filtered associative algebra whose associated graded is $\mathcal{A}_{G,R}$. Equivariant $K$-theory gives the multiplicative, $q$-difference version of the same Coulomb-branch construction.
\end{remark}

\subsection{Key examples}

\begin{example}[Pure gauge theory]
\label{ex:pure-gauge}
For $R = 0$ (no matter), $\mathcal{A}_{G,0} = \mathbb{C}[\mathfrak{t}^* \times T^\vee]^W$ where $T \subset G$ is a maximal torus, $T^\vee$ the dual torus, and $W$ the Weyl group. Thus $M_C(G, 0) = (\mathfrak{t}^* \times T^\vee)/W$, the ``universal centralizer'' moduli.
\end{example}

\begin{example}[Cotangent type]
\label{ex:cotangent}
For $G = GL(n)$ and $R = \mathbb{C}^n$ (the defining representation), $M_C \cong T^*(\mathbb{C}^n / S_n)$, a resolution of the symmetric product.
\end{example}

\begin{example}[Quiver Coulomb branches]
\label{ex:quiver-coulomb}
For a quiver gauge theory of type $A$, the Coulomb branch is a slice in the affine Grassmannian of $GL_N$, the \emph{affine Grassmannian slice} $\overline{\mathcal{W}}^\lambda_\mu$. These are isomorphic to (resolutions of) Nakajima quiver varieties for the \emph{dual} (Langlands dual, in many cases mirror) quiver. This is the mathematical incarnation of 3d mirror symmetry for type $A$ quivers.
\end{example}

%% =====================================================================
\section{$\CY_2$ quantum vertex chiral groups}
\label{sec:cy2-qvcg}
%% =====================================================================

\subsection{Quantum vertex chiral groups of hyperk\"ahler targets}

Since both $M_H(T)$ and $M_C(T)$ are $\CY_2$, the functor $\Phi$ of Theorem CY-A (Chapter~\ref*{ch:cy-to-chiral} of the main text) produces $\Etwo$-chiral algebras from their categories:
\[
 G(M_H) := \Phi(D^b(\Coh(M_H))), \qquad G(M_C) := \Phi(D^b(\Coh(M_C))).
\]
Each of these is a \emph{quantum vertex chiral group}: an $\Etwo$-chiral algebra whose representation category is braided monoidal, equipped with the root datum extracted from the $\CY_2$ geometry.

\begin{itemize}
 \item The $\mathbb{S}^2$-framing of $\HH_\bullet(D^b(\Coh(M)))$ provides the $\Etwo$-structure (the braiding).
 \item The CY trace $\Tr \colon \HH_2(\cC) \to k$ determines the modular characteristic $\kappa_{\mathrm{ch}}(G(M))$.
 \item The root datum $\mathcal{R}(M)$ is extracted from $K_0(M)$ with the Euler form and BPS multiplicities.
\end{itemize}

\subsection{Mirror symmetry for quantum vertex chiral groups}

The central conjecture:

\begin{conjecture}[Mirror symmetry for quantum vertex chiral groups]
\label{conj:mirror-qvcg}
Let $T$ and $T^!$ be 3d mirrors. Then the quantum vertex chiral groups of their branches are exchanged:
\begin{equation}
\label{eq:mirror-qvcg}
 G(M_H(T)) \cong G(M_C(T^!)), \qquad G(M_C(T)) \cong G(M_H(T^!)).
\end{equation}
More precisely, there is an equivalence of $\Etwo$-chiral algebras
\[
 \Phi(D^b(\Coh(M_H(T)))) \simeq \Phi(D^b(\Coh(M_C(T^!))))
\]
preserving the braided monoidal structure on representation categories, the modular characteristic, and the root datum.
\end{conjecture}

\begin{remark}
The conjecture has a tautological flavor if one simply applies $\Phi$ to both sides of $M_H(T) \cong M_C(T^!)$. The non-trivial content is:
\begin{enumerate}[label=(\alph*)]
 \item The \emph{algebraic structures} arising from the Higgs and Coulomb branch descriptions are very different: the Higgs side gives a CoHA (convolution on moduli of quiver representations), while the Coulomb side gives a BFN algebra (convolution on the affine Grassmannian). The conjecture asserts these produce equivalent $\Etwo$-chiral algebras.
 \item The \emph{root data} are a priori different: on the Higgs side, roots come from dimension vectors of quiver representations; on the Coulomb side, from cocharacters of the gauge group. The isomorphism of root data is itself a deep statement.
 \item The categorical form is stronger: $D^b(\Coh(M_H)) \simeq D^b(\Coh(M_C^!))$ as $\CY_2$ categories, not merely as an isomorphism of coarse varieties.
\end{enumerate}
\end{remark}

%% =====================================================================
\section{The CoHA of the Higgs branch}
\label{sec:coha-higgs}
%% =====================================================================

\subsection{Cohomological Hall algebras for quiver varieties}

For a quiver gauge theory with quiver $Q$ and potential $W$ (arising from the $\cN=4$ constraint), the Higgs branch is the Nakajima quiver variety $M_H = \mathfrak{M}(\mathbf{v}, \mathbf{w})$. The \emph{cohomological Hall algebra} of the Higgs branch is:
\[
 \CoHA(M_H) = \bigoplus_{\mathbf{d}} H^{\mathrm{BM}}_*(M_H^{(\mathbf{d})}, \mathbb{Q}),
\]
where $M_H^{(\mathbf{d})}$ denotes the moduli of sheaves/representations of dimension vector $\mathbf{d}$ on $M_H$, and the multiplication is the Hall-algebra convolution (Kontsevich--Soibelman, Schiffmann--Vasserot).

\begin{theorem}[CoHA = Yangian positive half]
\label{thm:coha-yangian}
For a quiver $Q$ without potential (i.e., the preprojective algebra case relevant to Nakajima quiver varieties):
\[
 \CoHA(Q) \cong Y^+(\frakg_Q),
\]
where $Y^+(\frakg_Q)$ is the positive half of the Yangian associated to the Kac--Moody algebra $\frakg_Q$ of the quiver.
\end{theorem}

This is the combined result of Schiffmann--Vasserot (for $\widehat{\mathfrak{gl}}_1$), Maulik--Okounkov (for general quivers via stable envelopes), and Davison (via critical CoHA techniques).

\subsection{$\Eone$-structure of the CoHA}

The CoHA is an \emph{associative} algebra; it is naturally an $\Eone$-algebra. The $\Etwo$ information is not carried by the Hall product alone. It enters through the stable-envelope $R$-matrix, equivalently through the braided representation category. In Volume~III terms:
\begin{itemize}
 \item $\CoHA(M_H) = Y^+(\frakg_Q)$ is the \emph{positive half} of the quantum vertex chiral group $G(M_H)$.
 \item It carries an $\Eone$-chiral structure: the ordered OPE from the Hall convolution.
 \item The full quantum vertex chiral group $G(M_H)$ has triangular decomposition
 \[
 G(M_H) = \frakn_+ \oplus \mathfrak{h} \oplus \frakn_-
 \]
 with $U(\frakn_+) = Y^+(\frakg_Q) = \CoHA(M_H)$.
 \item The passage from $\Eone$ (associative, the CoHA) to $\Etwo$ (braided representation theory) requires the $R$-matrix: this is the Maulik--Okounkov $R$-matrix, constructed via stable envelopes on Nakajima varieties.
\end{itemize}

\begin{remark}[$R$-matrix from stable envelopes]
\label{rem:stable-envelopes}
The Maulik--Okounkov stable envelope is a map
\[
 \mathrm{Stab}_{\mathfrak{C}} \colon H^*_T(M_H^T) \to H^*_T(M_H)
\]
depending on a choice of chamber $\mathfrak{C}$. The wall-crossing between chambers $\mathfrak{C}$ and $\mathfrak{C}'$ produces the $R$-matrix
\[
 R_{\mathfrak{C}, \mathfrak{C}'} = \mathrm{Stab}_{\mathfrak{C}'}^{-1} \circ \mathrm{Stab}_{\mathfrak{C}} \in \End(H^*_T(M_H^T)),
\]
satisfying the Yang--Baxter equation. This is the $\Etwo$-structure: the braiding of the quantum vertex chiral group $G(M_H)$.
\end{remark}

%% =====================================================================
\section{The Coulomb branch algebra as dual CoHA}
\label{sec:coulomb-dual-coha}
%% =====================================================================

\subsection{BFN algebra and the positive half of $G(M_C)$}

The BFN Coulomb branch algebra $\mathcal{A}_{G,R}$ (Definition~\ref{def:bfn}) is a commutative algebra whose spectrum is $M_C$. Its \emph{quantization} $\mathcal{A}_{G,R}^\hbar$ is an associative (non-commutative) algebra. The conjectural identification is:

\begin{conjecture}[BFN = positive half of $G(M_C)$]
\label{conj:bfn-positive}
The quantized BFN algebra $\mathcal{A}_{G,R}^\hbar$ is isomorphic to the positive half of the quantum vertex chiral group of the Coulomb branch:
\[
 \mathcal{A}_{G,R}^\hbar \cong Y^+(G(M_C)),
\]
where $Y^+(G(M_C))$ denotes the ``positive Yangian half'' of the quantum vertex chiral group $G(M_C)$.
\end{conjecture}

The evidence for this:
\begin{enumerate}[label=(\roman*)]
 \item For $G = GL(n)$, $R = N$ fundamentals, $M_C$ is an affine Grassmannian slice and $\mathcal{A}_{G,R}^\hbar$ is a truncated shifted Yangian (Kamnitzer--Webster--Weekes--Yacobi), which is literally a quotient of $Y^+(\frakg)$ for the appropriate $\frakg$.
 \item The BFN algebra has a convolution product coming from the affine Grassmannian, which is an $\Eone$-structure (associative, from the ``ordering'' of loop groups).
 \item The full quantum vertex chiral group $G(M_C)$ is expected from $\mathcal{A}_{G,R}^\hbar$ by the Drinfeld-double operation: add the negative half and the Cartan after constructing the required Hopf pairing and completion.
\end{enumerate}

\subsection{Mirror symmetry as Yangian duality}

Combining Theorem~\ref{thm:coha-yangian} and Conjecture~\ref{conj:bfn-positive}, 3d mirror symmetry at the level of quantum vertex chiral groups becomes:

\begin{equation}
\label{eq:yangian-duality}
 Y^+(\frakg_Q) = \CoHA(M_H(T)) \cong \mathcal{A}_{G^!,R^!}^\hbar = Y^+(G(M_C(T^!)))
\end{equation}

where $(G^!, R^!)$ is the gauge theory data of the mirror $T^!$. This is a ``Yangian-level mirror symmetry'': the CoHA of quiver representations on one side equals the quantized BFN algebra of the dual gauge theory on the other.

\begin{example}[Type $A$ mirror symmetry]
\label{ex:type-a-mirror}
For the $A_n$ quiver with dimension vector $\mathbf{v}$ and framing $\mathbf{w}$:
\begin{itemize}
 \item Higgs side: $M_H = \mathfrak{M}(\mathbf{v}, \mathbf{w})$, and $\CoHA(M_H) = Y^+(\widehat{\mathfrak{sl}}_n)$ (or a quotient/truncation depending on $\mathbf{v}, \mathbf{w}$).
 \item Coulomb side of mirror: $M_C(T^!)$ is an affine Grassmannian slice for $GL_N$ (where $N$ depends on $\mathbf{w}$), and $\mathcal{A}^\hbar_{T^!}$ is a truncated shifted Yangian for $\mathfrak{gl}_N$.
 \item The isomorphism $Y^+(\widehat{\mathfrak{sl}}_n) \cong$ (truncated shifted Yangian for $\mathfrak{gl}_N$) is a theorem in the mathematics literature (Kamnitzer et al.), and it \emph{is} 3d mirror symmetry at the algebra level.
\end{itemize}
\end{example}

%% =====================================================================
\section{Higgs bundles from the Higgs branch}
\label{sec:higgs-bundles}
%% =====================================================================

\subsection{Pure gauge theory on a curve}

Fix a smooth projective curve $C$ and a reductive group $G$. Consider the 3d $\cN=4$ pure gauge theory (no matter, $R = 0$) ``compactified on $C$,'' or more precisely, consider the dimensional reduction of 4d $\cN=2$ pure gauge theory on $C$. In this setup:

\begin{theorem}[Higgs branch = Hitchin moduli]
\label{thm:higgs-hitchin}
The Higgs branch of the 3d $\cN=4$ theory associated to $G$ on $C$ is the moduli space of Higgs bundles:
\[
 M_H = T^*\Bun_G(C) = \Higgs(C, G),
\]
the cotangent bundle to the moduli stack of principal $G$-bundles on $C$. Concretely,
\[
 \Higgs(C, G) = \left\{(P, \varphi) \;\middle|\; P \in \Bun_G(C),\; \varphi \in H^0(C, \mathrm{ad}(P) \otimes K_C)\right\},
\]
where $\varphi$ is the Higgs field.
\end{theorem}

\begin{proof}[Sketch]
The Higgs branch is $T^*R \mathbin{/\!\!/\!\!/} G$ with $R$ the infinite-dimensional space of connections on $C$. The hyperk\"ahler quotient gives Hitchin's equations, whose complex-structure-$I$ description is exactly the moduli of Higgs bundles.
\end{proof}

\begin{remark}
\label{rem:higgs-cy2}
The Hitchin moduli space $\Higgs(C, G)$ is hyperk\"ahler (Hitchin, 1987) and therefore $\CY_2$ (more precisely, it is a non-compact hyperk\"ahler manifold; its derived category is a $\CY_2$ category). The Hitchin fibration $h \colon \Higgs(C,G) \to \mathcal{B}$ (the Hitchin base, a vector space of invariant polynomials) is a completely integrable system.
\end{remark}

\subsection{The quantum vertex chiral group of Higgs bundles}

Since $\Higgs(C, G)$ is $\CY_2$, it carries a quantum vertex chiral group:
\[
 G(\Higgs(C,G)) = \Phi(D^b(\Coh(\Higgs(C,G)))).
\]

\begin{conjecture}[Hitchin quantum vertex chiral group]
\label{conj:hitchin-qvcg}
The quantum vertex chiral group $G(\Higgs(C,G))$ satisfies:
\begin{enumerate}[label=(\roman*)]
 \item Its positive half $Y^+(G(\Higgs(C,G)))$ is the CoHA of Higgs sheaves on $C$, the cohomological Hall algebra of the moduli of Higgs sheaves on $C$ studied by Minets, Sala--Schiffmann, and Kapranov--Vasserot.

 \item The full quantum vertex chiral group is the ``Higgs Yangian'' or ``Hitchin Yangian'' $Y(\Higgs(C,G))$, an algebra that organizes the BPS spectrum of the Hitchin system.

 \item The modular characteristic $\kappa_{\mathrm{ch}}(G(\Higgs(C,G)))$ is related to the Euler characteristic of $\Higgs(C,G)$, which is computed by the topological field theory of Hausel--Thaddeus and Hausel--Rodriguez-Villegas and (for $G = GL_n$) is intimately related to the $P = W$ conjecture (now theorem, by Maulik--Shen--Yin and Hausel--Mellit--Minets--Schiffmann).
\end{enumerate}
\end{conjecture}

\subsection{Mirror symmetry for Hitchin moduli}

By 3d mirror symmetry, the Coulomb branch of the theory $(G, C)$ is related to the Higgs branch of the mirror. In the Hitchin context, this takes the following form:

\begin{conjecture}[Hitchin mirror symmetry via quantum vertex chiral groups]
\label{conj:hitchin-mirror}
For the 3d $\cN=4$ theory associated to $(G, C)$:
\begin{enumerate}[label=(\roman*)]
 \item The Coulomb branch is $M_C = \Higgs(C, G^\vee)$, where $G^\vee$ is the Langlands dual group.

 \item The mirror symmetry $M_H(G, C) \cong M_C(G^\vee, C)$, i.e.,
 \[
 \Higgs(C, G) \cong \Higgs(C, G^\vee)
 \]
 (as hyperk\"ahler manifolds, with a rotation of the hyperk\"ahler structure), is the geometric content of the Strominger--Yau--Zaslow proposal for Hitchin moduli, proved (at the level of generic Hitchin fibers) by Hausel--Thaddeus.

 \item At the quantum vertex chiral group level:
 \[
 G(\Higgs(C, G)) \cong G(\Higgs(C, G^\vee)),
 \]
 the ``CoHA of $G$-Higgs sheaves on $C$'' is equivalent to the ``quantized BFN algebra of $G^\vee$ on $C$.''
\end{enumerate}
\end{conjecture}

\begin{remark}[Connection to geometric Langlands]
\label{rem:geom-langlands}
This is 3d mirror symmetry as the \emph{physical origin of geometric Langlands}. The program of Kapustin--Witten (2006) derives the geometric Langlands correspondence from the compactification of 4d $\cN=4$ super Yang--Mills on $C$:
\begin{itemize}
 \item The A-model on $\Higgs(C, G)$ (with complex structure $J$) = the B-model on $\Higgs(C, G^\vee)$ (with complex structure $I$).
 \item The category of A-branes on $\Higgs(C, G)$ = $D$-modules on $\Bun_{G^\vee}(C)$ = geometric Langlands dual.
 \item In our language: the quantum vertex chiral group $G(\Higgs(C, G))$ in complex structure $J$ is equivalent to $G(\Higgs(C, G^\vee))$ in complex structure $I$. This is a \emph{chiral-algebraic refinement} of the Kapustin--Witten story.
\end{itemize}
This provides the bridge to Chapter~\ref*{ch:geometric-langlands} of the main text: the quantum vertex chiral groups of Hitchin moduli \emph{are} the algebraic structure underlying the geometric Langlands correspondence.
\end{remark}

%% =====================================================================
\section{The big picture: a dictionary}
\label{sec:dictionary}
%% =====================================================================

We summarize the identifications developed in this note.

\medskip

\begin{center}
\renewcommand{\arraystretch}{1.5}
\begin{tabular}{@{}p{3.8cm}p{4.2cm}p{4.5cm}@{}}
\hline
\textbf{3d $\cN=4$ physics} & \textbf{$\CY_2$ geometry} & \textbf{Quantum vertex chiral group} \\
\hline
Higgs branch $M_H$ & Nakajima quiver variety $\mathfrak{M}(\mathbf{v}, \mathbf{w})$ & Root datum from quiver $Q$ \\
Coulomb branch $M_C$ & BFN variety $\Spec(\mathcal{A}_{G,R})$ & Root datum from cocharacters of $G$ \\
3d mirror symmetry $M_H(T) \cong M_C(T^!)$ & Hyperk\"ahler isomorphism & $G(M_H(T)) \cong G(M_C(T^!))$ \\
CoHA of $M_H$ & Hall convolution on $\mathfrak{M}$ & $Y^+(\frakg_Q)$ = positive half \\
Quantized BFN algebra $\mathcal{A}^\hbar_{G,R}$ & Affine Grassmannian convolution & Positive half of $G(M_C)$ \\
Maulik--Okounkov $R$-matrix & Stable envelopes & $\Etwo$-structure (braiding) \\
Pure gauge on $C$: $M_H = \Higgs(C,G)$ & Hitchin moduli space & $G(\Higgs(C,G))$ \\
$M_H(G) \cong M_C(G^\vee)$ & SYZ for Hitchin fibers & $G(\Higgs(C,G)) \cong G(\Higgs(C,G^\vee))$ \\
Kapustin--Witten & A-model/B-model exchange & Geometric Langlands via $\Etwo$-chiral \\
\hline
\end{tabular}
\end{center}

\medskip

%% =====================================================================
\section{Structural analysis and open problems}
\label{sec:open}
%% =====================================================================

\subsection{The $\CY_2$ vs $\CY_3$ distinction}

The toric $\CY_3$ story of Chapter~\ref*{ch:toric-coha} in the main text and the $\CY_2$ story of this note are complementary but structurally different:

\begin{itemize}
 \item \textbf{$\CY_3$}: The target $X$ is a Calabi--Yau threefold. The CoHA arises from curve-counting (DT invariants) on $X$ itself. The specialised output $A_X$ is natively $\Eone$-chiral on the reference curve. The non-symmetric $\Etwo$ quantum-group structure belongs to $\mathcal{Z}(\Rep^{\Eone}(A_X))$ or to the corresponding Drinfeld double once the Hopf pairing, completion, and half-braiding data have been constructed. The root multiplicities come from BPS invariants of $X$.

 \item \textbf{$\CY_2$ from 3d $\cN=4$}: The target $M$ is a hyperk\"ahler surface (the branch of a 3d gauge theory). The CoHA arises from sheaf counting on $M$ or equivalently from quiver representation theory. The quantum vertex chiral group $G(M)$ has $\Etwo$-structure from the $\mathbb{S}^2$-framing of $\HH_\bullet(D^b(\Coh(M)))$. The root multiplicities come from dimensions of quiver representation moduli.
\end{itemize}

The $\CY_2$ case is simpler (lower dimension, explicit hyperk\"ahler metric available in many cases) but the mirror symmetry structure is richer because one has \emph{two} branches to compare.

\subsection{Quantization levels}

There are several layers of quantization to track:

\begin{enumerate}[label=(\arabic*)]
 \item \textbf{Classical}: $M_H$ and $M_C$ as Poisson varieties. Mirror symmetry is a Poisson isomorphism.
 \item \textbf{Deformation quantization}: the quantized Coulomb branch algebra $\mathcal{A}^\hbar_{G,R}$ and the enveloping algebra of the CoHA. Mirror symmetry is an algebra isomorphism.
 \item \textbf{Categorical}: $D^b(\Coh(M_H)) \simeq D^b(\Coh(M_C^!))$ as $\CY_2$ categories. Mirror symmetry is a CY equivalence.
 \item \textbf{Chiral}: $G(M_H(T)) \cong G(M_C(T^!))$ as $\Etwo$-chiral algebras. Mirror symmetry is an equivalence of quantum vertex chiral groups.
\end{enumerate}

Each level refines the previous, and the chiral level (4) is the most refined statement, incorporating all the modular/genus data from the shadow obstruction tower.

\subsection{Open problems}

\begin{enumerate}[label=\textbf{P\arabic*.}]
 \item \textbf{Rigorous CoHA = BFN.} Prove that for a 3d mirror pair $(T, T^!)$, the CoHA of the Higgs branch of $T$ is isomorphic (as an $\Eone$-algebra) to the quantized BFN algebra of $T^!$. This is known for type $A$ quivers (Kamnitzer--Webster--Weekes--Yacobi); the general case is open.

 \item \textbf{$R$-matrix matching.} Show that the Maulik--Okounkov $R$-matrix on the Higgs side matches the $R$-matrix arising from the quantized BFN algebra on the Coulomb side. This establishes the $\Etwo$ equivalence after the $\Eone$ algebra equivalence.

 \item \textbf{Modular characteristic under mirror symmetry.} Compute $\kappa_{\mathrm{ch}}(G(M_H))$ and $\kappa_{\mathrm{ch}}(G(M_C^!))$ independently and verify they agree. The Higgs-branch modular characteristic is a concrete extraction problem from the Hausel--Rodriguez-Villegas formula for the mixed Hodge polynomial of quiver varieties.

 \item \textbf{Hitchin CoHA.} Give a complete construction of the CoHA of Higgs sheaves on a curve $C$ for general $G$ and identify it with the positive half of a quantum vertex chiral group. The work of Sala--Schiffmann and Minets handles the $GL_n$ case; the general case requires new ideas.

 \item \textbf{Shadow obstruction tower from BFN.} Construct the full shadow obstruction tower $\Theta^{\leq r}$ (in the sense of Volume~I) directly from the BFN data. The degree-2 component $\kappa_{\mathrm{ch}}$ is the candidate BFN Weyl vector; higher degrees encode the expected quantum corrections to the Coulomb branch metric.

 \item \textbf{Wall-crossing.} Interpret the Kontsevich--Soibelman wall-crossing formula for DT invariants of $\CY_2$ categories as a gauge transformation of the MC element $\Theta_{G(M)}$. For the Hitchin system, the target comparison is the wall-crossing of spectral networks (Gaiotto--Moore--Neitzke).

 \item \textbf{Symplectic duality.} The mathematical formalization of 3d mirror symmetry is Braden--Licata--Proudfoot--Webster's ``symplectic duality.'' Formulate the quantum vertex chiral group construction in their framework and verify Conjecture~\ref{conj:mirror-qvcg} as a consequence of symplectic duality.
\end{enumerate}

%% =====================================================================
\section{Conclusion}
\label{sec:conclusion}
%% =====================================================================

The message of this note is that 3d $\cN=4$ mirror symmetry, when viewed through the lens of $\CY_2$ quantum vertex chiral groups, provides:

\begin{enumerate}[label=(\roman*)]
 \item A \emph{physical origin} for the quantum vertex chiral groups $G(M)$ of hyperk\"ahler targets: they arise as the algebraic structure organizing the BPS spectrum of 3d $\cN=4$ theories.

 \item A \emph{duality} on quantum vertex chiral groups: 3d mirror symmetry exchanges $G(M_H(T))$ and $G(M_C(T^!))$, providing a non-trivial isomorphism between CoHA-type and BFN-type constructions.

 \item A \emph{direct connection to geometric Langlands}: the quantum vertex chiral group of the Hitchin moduli space $\Higgs(C, G)$ is the chiral-algebraic structure underlying the Kapustin--Witten approach to geometric Langlands, and Langlands duality $G \leftrightarrow G^\vee$ is 3d mirror symmetry for this class of theories.

 \item A \emph{bridge between $\CY_2$ and $\CY_3$}: the Higgs and Coulomb branches are $\CY_2$, but the 3d theory itself can arise from compactification of a theory on a $\CY_3$ (or from engineering on a $\CY_3$ singularity). The quantum vertex chiral group construction provides a uniform framework encompassing both dimensions.
\end{enumerate}

The Coulomb-branch and Hitchin-moduli constructions supply the $\CY_2$ mirror-symmetric input for the braided factorisation and geometric-Langlands parts of the main text.

\end{document}
